\documentclass[11pt,a4paper]{article}
\usepackage[T1]{fontenc}
\usepackage[utf8]{inputenc}
\usepackage[french]{babel}
\frenchsetup{SmallCapsFigTabCaptions=false}
% The English language files of this Babel installation are broken, so the
% French language is loaded for hyphenation only and every caption is renamed.
\addto\captionsfrench{%
  \renewcommand{\contentsname}{Contents}%
  \renewcommand{\abstractname}{Abstract}%
  \renewcommand{\figurename}{Figure}%
  \renewcommand{\tablename}{Table}%
  \renewcommand{\refname}{References}%
}
\usepackage{lmodern,microtype,amsmath}
\usepackage[margin=23mm]{geometry}
\usepackage{booktabs,array,tabularx,longtable}
\usepackage{xcolor,tikz,enumitem,float}
\usepackage{hyperref}
\usetikzlibrary{arrows.meta,positioning,fit,calc}
\definecolor{ink}{HTML}{1F242C}
\definecolor{muted}{HTML}{5C626C}
\definecolor{line}{HTML}{D9DEE5}
\definecolor{accent}{HTML}{AF0032}
\definecolor{rose}{HTML}{FDF0F4}
\hypersetup{colorlinks=true,linkcolor=black,citecolor=black,urlcolor=ink,
 pdftitle={LANCE: architecture, protocols and current state},
 pdfauthor={LANCE project},pdflang={en-GB}}
\urlstyle{same}
\def\UrlBreaks{\do\/\do\-\do\_\do\.\do\:}
\newcolumntype{Y}{>{\raggedright\arraybackslash}X}
\DeclareRobustCommand{\code}[1]{{\small\path{#1}}}
\newcommand{\source}[2]{\href{../../#1}{\nolinkurl{#1}}. #2}
\newcommand{\revision}{2026-10-30 15:45 (Europe/Brussels, UTC+01:00)}
\setlength{\parskip}{3pt}
\setlength{\emergencystretch}{2em}
\setlist{nosep,leftmargin=5mm,topsep=3pt}
\clubpenalty=10000 \widowpenalty=10000
\setcounter{tocdepth}{1}

\title{LANCE\\[4pt]
\large Architecture, protocols and current state\\[2pt]
\normalsize Report as of 30 October 2026}
\author{LANCE project\\\small Written by specialist strands, independently reviewed and synthesised}
\date{Last revision: \revision\\[3pt]
\small Documentation state as of 30 October 2026; no new experiment is reported}

\begin{document}
\maketitle

\begin{abstract}
LANCE is a platform for the security audit of authorised connected-object
networks, developed in the context of the UMONS work of the NATO SPS G6259
programme. This report describes the pipeline architecture, the audit protocols
and their evidence contracts, the experimentation harnesses and their hermetic
sandbox, the execution infrastructure and the delivery pipeline, and the actual
state of the benchmark. It succeeds the 15 October 2026 overview, which skimmed
over the architecture and the protocols: those parts are developed here down to
the code contract.

The document consistently distinguishes what is implemented, what is measured and
what remains proposed. No new experiment is reported: the measurements cited are
those already recorded, with their date. The limits of each component are
explicit, including those of the sandbox, whose guard is not an exhaustive
native-code sandbox.
\end{abstract}

\tableofcontents
\clearpage

\section{Executive summary}

This page answers four questions without jargon: what the system does, where it
stands, what is proven, and what has to be decided.

\subsection{What the system does}

LANCE conducts a security audit on a network of connected objects \emph{whose
ownership and authorisation are established}. It explores the network, issues
hypotheses of flaws, attempts to verify them by observable evidence, and then
writes a report. Two policies can carry out this audit: a rule-based automation
and an agent based on a language model. The project seeks to determine which of
the two succeeds better, with identical tools, evidence and budgets.

The audit runs in six phases, each producing a deliverable the next one consumes:
the topology graph, the reconnaissance, the vulnerability hypotheses, the
exploitation attempts, the intrusion campaign, and the report. Every tool call
passes through a single boundary that archives its arguments, its result and its
cost in a journal; that journal, not the policy's own conclusion, is what the
evaluator recounts.

A run can use one of four execution modes. The full pipeline runs all six phases
with a model. A rule-based mode replaces the model with declared rules for
analysis and verification, on phases three and four only. A bare mode gives the
model a bounded tool loop on an address range, without scenario or score. A
scripted mode runs versioned deterministic procedures. Of the four, only the full
pipeline is scored against a ground truth; the others produce results that are
read but not benchmarked.

The system runs in a laboratory: scenarios are deployed on a Proxmox host, audited
over a dedicated network, then torn down. A browser interface, served by three
instances, starts runs, streams their events and reads their results.

\subsection{Where the project stands}

\begin{table}[H]
\centering
\caption{Status in one line per strand.}
\label{tab:etat-volet}
\small
\begin{tabular}{@{}p{3.3cm}p{1.9cm}p{8.3cm}@{}}
\toprule
Strand & Status & Detail \\
\midrule
Audit engine & mature & Six phases in service, four modes, execution boundary that logs every call \\
Evidence & mature & Briefing/oracle/evaluator contract applied to five families; the log is authoritative, not the decision \\
Harness & consolidated & Core shared since 29 October; sandbox executable outside the laboratory \\
Infrastructure & restored & Three instances, automatic delivery pipeline, seven defects fixed \\
Benchmark & partial & One hypothesis measured out of six, one family rejected, no inference computable \\
Scientific proof & absent & No superiority of the model has been established to date \\
\bottomrule
\end{tabular}
\end{table}

\subsection{What is proven, and what is not}

\textbf{Proven.} Every tool call is traced and bounded; a lost piece of evidence
interrupts the run. A harness can run outside the laboratory, under a guard whose
blind spots are declared. On the best-documented family, a third script, which
reads by convention and uses no model, obtains the same result as the agent,
which \emph{absorbs} the gain attributed to the model. The corpus is mostly
detectable by signature, which makes a null gap expected by construction.

\textbf{Not proven.} That a model does better than a competent script. That a gap
is significant. That the results generalise beyond eight cases and one property.

\subsection{The three decisions that govern what follows}

\begin{enumerate}
\item \textbf{The independent families.} The inference tool was repaired on
30 October 2026 and now refuses to produce a number it cannot defend. What
remains is the families themselves: the corpus holds at most three defensible
ones, and a decision needs five. Either build them, applying the filter that one
hour of observation established, or publish a bounded methodological result and
declare it as such.
\item \textbf{The opening condition.} Several families remain undecidable for lack
of a written reference automation. Either write it, or explicitly renounce these
families.

\item \textbf{The comparable arm.} A comparison of the two policies supposes that
both receive the same observations. Four mechanisms currently favour the rules:
the phase-four probes only they may emit, a supplement that rewrites the scan
files, the scanner findings added to every aggregation, and a benchmark score
reserved for the full pipeline. Either give the model an equivalent access, or
state that the comparison measures verification rather than interpretation.
\end{enumerate}

\subsection{How to read this report}

Section~\ref{sec:architecture-pipeline} describes the path of an audit.
Section~\ref{sec:protocoles} describes what makes a result verifiable.
Section~\ref{sec:harnais-sandbox} describes the experimental tooling.
Section~\ref{sec:infrastructure} describes the machine and the delivery.
Section~\ref{sec:benchmark} establishes what is measured.
The conclusion confronts these areas and notes their contradictions; it does not
repeat the summary above.

Each statement refers to the file that supports it. Behaviours that the code does
not allow to be established are flagged as such rather than filled in.

\clearpage
\section{System overview}
\label{sec:vue-ensemble}

Before going into the contracts, a map of the system and of the laboratory. It
situates the components described below and shows where evidence passes.

\begin{figure}[H]
\centering
\begin{tikzpicture}[
  node distance=8mm and 9mm,
  box/.style={draw=line, fill=white, rounded corners=2pt, line width=.55pt,
    text=ink, align=center, inner sep=5pt, font=\sffamily\footnotesize,
    text width=29mm, minimum height=12mm},
  focal/.style={box, draw=accent, fill=rose, line width=.75pt},
  flow/.style={-{Latex[length=1.8mm,width=1.1mm]}, draw=muted, line width=.6pt}
]
\node[box] (lab) {Simulated IoT\\laboratory};
\node[box, right=of lab] (run) {Audit engine\\six phases};
\node[focal, right=of run] (tools) {Execution boundary\\evidence log};
\node[box, below=of lab] (eval) {Evaluator\\offline};
\node[box, right=of eval] (verdict) {Verdict};
\node[box, above=of eval] (art) {Run\\artefacts, 01 to 06};
\node[box, above=of verdict] (oracle) {Private oracle\\fixed before the run};

\draw[flow] (lab) -- (run);
\draw[flow] (run) -- (tools);
\draw[flow] (run) -- (art);
\draw[flow] (art) -- (eval);
\draw[flow] (oracle) -- (verdict);
\draw[flow] (eval) -- (verdict);
\draw[flow, dashed] (tools.south) -- (verdict.north east);
\end{tikzpicture}
\caption{The journey of a piece of evidence. The tools produce logged
observations; the private oracle, written before execution, makes it possible to
recompute the verdict from the log without believing the system's decision.}
\label{fig:tripwire}
\end{figure}

Three properties can be read from this map.

\textbf{The log is the source.} A result is not the conclusion written by the
policy, but what the evaluator recomputes from the archived observations. A cited
piece of evidence that does not appear in the log is rejected.

\textbf{The oracle precedes the run.} It is fixed before the policy runs, and a
finding that contradicts it invalidates the trial instead of correcting it. This
precedence is what prevents the answer from being adjusted after the fact.

\textbf{The sandbox is not the measurement.} It executes the deterministic arms
without a model or a laboratory; it instruments the contrast, it does not measure
it.

\begin{figure}[H]
\centering
\begin{tikzpicture}[
  node distance=3mm and 5mm,
  h/.style={draw=line, fill=white, rounded corners=2pt, line width=.55pt,
    text=ink, align=center, inner sep=3pt, font=\sffamily\footnotesize},
  dev/.style={h, fill=rose, draw=accent},
  flow/.style={-{Latex[length=1.6mm,width=1mm]}, draw=muted, line width=.6pt}
]
\node[h, text width=22mm] (poste) {Development\\workstation};
\node[h, text width=24mm, right=of poste] (ci) {GitHub Actions\\validation};
\node[dev, text width=26mm, right=of ci] (lxc) {Master container\\three instances};
\node[h, text width=22mm, right=of lxc] (umons) {UMONS relay\\model};

\draw[flow] (poste) -- node[above, font=\sffamily\tiny, text=muted] {push} (ci);
\draw[flow] (ci) -- node[above, font=\sffamily\tiny, text=muted] {deploys} (lxc);
\draw[flow] (lxc) -- node[above, font=\sffamily\tiny, text=muted] {calls} (umons);
\end{tikzpicture}
\caption{The delivery pipeline. No execution work takes place on the development
workstation: the push triggers validation, which deploys to the container.}
\label{fig:circuit}
\end{figure}

\clearpage


\section{Scope and method}
\label{sec:scope}

This report answers a reading question: \emph{how is the system built, and where
does it actually stand?} It does not replace the reference pages of \code{docs/},
which describe the current behaviour, nor the detailed protocols of
\code{research/}, which carry the reasoning.

Each part was written by a specialist working on the \code{dev/1} branch, then
reviewed against that same code by an independent reader. A cross-cutting
synthesis was produced separately. Statements that cannot be verified in the
repository are flagged as such rather than filled in.

The report is dated: it describes a state, not a contract. Later behaviour may
contradict it without making it false.

% === PARTIES RÉDIGÉES PAR LES AGENTS ===
\section{End-to-end pipeline architecture}
\label{sec:architecture-pipeline}

This part describes the effective path of a LANCE audit, from the selection of an
execution mode to the files deposited in the run folder, relying exclusively on
the code of the \code{dev/1} branch. The statements refer to the module and,
where useful, to the function that implements them; behaviours that the code does
not allow to be established are flagged as such rather than filled in by
inference.

\subsection{Overview and mode selection}

The application does not carry the audit logic: it validates the request,
guarantees a single active run per instance, streams the SSE events and re-reads
the artefacts (\code{docs/architecture/run-modes.md}). It is \code{make\_runner}
(\code{src/agent/run\_modes.py}) that chooses the execution strategy. For
\code{lance} and \code{rules}, it instantiates a \code{Pipeline} and wraps it in
\code{LanceRunner}, an adapter that reuses the pipeline object and its public
API. For \code{vanilla} and \code{scripted}, it builds a \code{SimpleRunner}, a
standalone class that shares with the pipeline only the tool execution boundary,
the evidence log, the cost tracking and the laboratory reservation.

\begin{table}[H]
\centering
\caption{Execution modes and corresponding boundary.}
\label{tab:run-modes}
\small
\begin{tabular}{@{}llll@{}}
\toprule
Mode & Engine & Model & Effects \\
\midrule
\code{lance} & \code{Pipeline} (6 phases) & required & scenario, deployment, evaluation \\
\code{rules} & \code{Pipeline} (phases 3--4) & none & scenario and deterministic audit \\
\code{vanilla} & \code{SimpleRunner} & required & bounded Nmap probe on CIDR \\
\code{scripted} & \code{SimpleRunner} & none & versioned procedures \\
\bottomrule
\end{tabular}
\end{table}

The \code{rules} mode is not a second engine: \code{make\_runner} sets
\code{decision\_policy="rules"}, \code{phases=[3, 4]} and
\code{execution\_profile="full"}, and rejects any provider, direct network or
monetary budget. The \code{vanilla} and \code{scripted} modes, conversely, reject
\code{scenario\_id} and any harness-specific setting. The qualifier
\code{full}/\code{compact} designates the orchestration of the pipeline, never the
type of runner.

\subsection{The six phases and their deliverables}

The declarative registry \code{src/agent/registry.py} (\code{AGENTS}) assigns to
each phase a prompt template, a tool group, a validator and a deliverable name.
The routing is explicit in \code{src/agent/phases/registry.py}:
\code{PHASE\_MODULES} maps the numbers 1 to 6 to fixed modules, without dynamic
import.

\begin{table}[H]
\centering
\caption{Phases, modules, main deliverables and associated artefacts.}
\label{tab:phases}
\small
\begin{tabular}{@{}lp{2.7cm}p{3.2cm}p{4.05cm}@{}}
\toprule
Phase & Module & Main deliverable & Associated artefacts \\
\midrule
1 & \code{graph/run.py} & \code{01\_graph\_analysis.md} & \code{01\_graph\_evidence.json} \\
2 & \code{recon/run.py} & \code{02\_recon.md} & \code{02\_recon\_evidence.json}, \code{02\_topology\_edges.json} \\
3 & \code{analysis/run.py} & \code{03\_vuln\_analysis.json} & \code{03\_vuln\_analysis\_raw.json}, \code{03\_phase3\_status.json}, \code{03\_device\_*} \\
4 & \code{verification/run.py} & \code{04\_exploitation.json} & \code{execution\_errors} diagnostics \\
5 & \code{intrusion/run.py} & \code{05\_intrusion.json} & \code{05\_intrusion\_context.json} \\
6 & \code{report/run.py} & \code{06\_report.md} & \code{06\_report\_analysis.md}, per-section sheets \\
\bottomrule
\end{tabular}
\end{table}

Phase 1 projects the results of the graph tools into a log of facts independent
of the model's discourse
(\code{GraphPhase.\_build\_graph\_evidence\_projection}), and the runner rejects
a submission whose per-device coverage is incomplete
(\code{src/agent/core/runner.py}). Phase 2 lets the model choose the probes, but
\code{02\_recon.md} is rendered from the tool log, not copied by the model. Phase
3 runs in three steps: \code{scan.scan\_phase} (deterministic observations),
\code{devices.analyze\_devices} (analysts per device) and then
\code{run.aggregate\_phase}, which normalises, deduplicates and projects the
findings and writes \code{03\_vuln\_analysis.json}. Phase 4 launches micro-agents
per vulnerability and aggregates their results. Phase 5 produces its assessment
from the log: the model supplies only an end marker, and the status remains
\code{objectives\_status=not\_certified}. Phase 6 is particular: its
\code{AgentConfig} declares no tool and \code{max\_turns=1}; the report is
assembled deterministically from sheets, the model writing only a bounded analyst
note (\code{06\_report\_analysis.md}).

\subsection{Graph and run context}

The graph is built in two layers. \code{src/loader.py} reads the infrastructure
YAML into dataclasses (\code{load\_yaml}) and then populates a backend
(\code{build\_graph}, default backend \code{NetworkXBackend}).
\code{src/graph\_backend.py} defines the abstract interface \code{GraphBackend}
(\code{get\_device}, \code{get\_neighbors}, \code{find\_all\_paths},
\code{get\_attack\_surface}, \code{get\_graph\_stats}, \code{to\_dict}) and its
NetworkX implementation on a \code{DiGraph}. Two details deserve to be stressed:
\code{get\_neighbors} returns the sorted union of successors and predecessors, and
\code{find\_all\_paths} works on an undirected view (\code{to\_undirected}) in
order to find paths in both directions, including through the abstract backend.

The state belongs to the run. Each \code{Pipeline} creates an immutable
\code{RunContext} (run folder and \code{GraphContext}); \code{GraphContext}
carries the backend, the infrastructure, the CVE reports, the risk scores, the
scenario or discovery topology, the weighted graph and the latest delta. The graph
tools retrieve this context via a \code{ContextVar}
(\code{src/agent/core/graph\_context.py}), without modifying the arguments exposed
to the model; \code{graph\_bound} and \code{bind\_graph\_callable} explicitly
capture the owner, including for thread workers, rather than relying on implicit
context propagation.

\subsection{The tool governance boundary}

All profiles and all modes go through \code{wrap\_tool}
(\code{src/agent/core/executor.py}). On each call, the wrapped function:
(1) checks that the previous piece of evidence could be archived, cooperative
stop and the budgets via \code{check\_execution\_limits}; (2) atomically
increments a counter and refuses the call if \code{max\_tool\_calls} is reached;
(3) assigns a reference \code{tc-<uuid>}; (4) applies the scope refusals;
(5) executes the tool; (6) systematically archives an entry in
\code{tool\_calls.jsonl}, including for a refusal or an exception.

\begin{table}[H]
\centering
\caption{Controls applied by \code{check\_execution\_limits} and \code{wrap\_tool}.}
\label{tab:executor}
\small
\begin{tabular}{@{}p{2.6cm}p{4.6cm}p{6.0cm}@{}}
\toprule
Control & Mechanism & Consequence of a failure \\
\midrule
Evidence integrity & flag \code{\_evidence\_integrity\_failed} & \code{EvidenceWriteError}, run not verifiable \\
User stop & \code{\_stop\_event} (experimental scope) & \code{RunStopped} \\
Time budget & \code{max\_duration\_s} via \code{time.monotonic()} & \code{BudgetExceeded} \\
Monetary budget & \code{tracker.check\_budget()} & \code{BudgetExceeded} \\
Call budget & counter under \code{\_artifact\_log\_lock} & \code{BudgetExceeded} \\
Scope & \code{\_intrusion\_scope\_violation} & refused call archived as JSON \\
Free execution & \code{python\_exec} & refusal \code{unverifiable\_execution\_scope} \\
\bottomrule
\end{tabular}
\end{table}

Two properties are important for the interpretation of the results. On the one
hand, a failure to write the log is never ignored: it is converted into
\code{EvidenceWriteError}, which interrupts the run rather than leaving untraced
actions. On the other hand, the cost limit of \code{check\_budget} applies to
observed usage (\code{total\_cost()}), not to a reservation of in-flight calls;
the code docstring says so explicitly. The scope is checked on the destination
fields (\code{ip}, \code{host}, \code{target}, \code{url}) against the
\code{target\_subnet} of the run; an action without a verifiable destination is
refused upstream (\code{python\_exec}), and the module header specifies that these
are application guards, not a system sandbox.

\subsection{Artefacts and cost tracking}

The run folder is created exclusively under \code{output/agent/}
(\code{Pipeline.\_\_init\_\_}); in the event of a collision, a hexadecimal suffix
is appended and no earlier folder is overwritten. \code{run\_meta.json} records
phase, profile, policies, commit and limits; \code{cost\_summary.json} is written
in the \code{finally} block of the pipeline, and therefore even on exception. The
collection contract (\code{src/benchmark/artifacts.py}) defines a positive list of
allowed files (\code{\_ALLOWED\_ROOT\_FILES}) and \code{03\_device\_*.json}
patterns; \code{create\_submission\_bundle} produces a zip equipped with a
SHA-256 manifest and re-checks the manifest after writing. On the agent side, a
deliverable attempt is archived in \code{.attempts/} and recorded in
\code{deliverable\_attempts.jsonl}; the route
\code{GET /api/runs/\{id\}/download/zip} excludes private files and hidden
entries.

Cost tracking (\code{src/agent/cost\_tracker.py}) resolves the price once per
phase (\code{\_resolve\_pricing}) and records it in \code{PhaseUsage}, so that a
completed run cannot be re-evaluated after the fact. Each phase carries the source
of the price (\code{dynamic\_catalog}, \code{static\_catalog},
\code{subscription}, \code{default\_estimate}) and a \code{cost\_is\_estimate}
indicator. \code{summary()} explicitly distinguishes the cumulative duration of
the agents (\code{total\_agent\_duration\_s}) from the wall-clock time
(\code{wall\_clock\_duration\_s}); the former does not measure the time actually
elapsed.

\subsection{What the architecture does not guarantee}

The architecture is isolated per run for the deliverables and the graph, but it is
not entirely so: \code{src/agent/phases/README.md} indicates that the phase
services may still call the common engine and that the simultaneity of scenarios
on a shared laboratory is not guaranteed. \code{reserve\_lab} is only active if
\code{LANCE\_LAB\_LOCK} is set and is only reentrant in the same thread. Moreover,
without an explicit parameter, the interactive pipeline imposes neither a call
ceiling nor an analysis timeout per device: the budgets of
table~\ref{tab:executor} are therefore conditional.

Three points could not be established by reading the code. \emph{One}: no
mechanism for archiving or compacting run folders was found in \code{src/} (no
occurrence of \code{tarfile}, \code{shutil.make\_archive} or equivalent);
retention appears to be that of the file system, which the documentation does not
describe. \emph{Two}: \code{02\_recon\_evidence.json} is read by
\code{src/agent/scanner.py} and \code{src/agent/core/runner.py}, but its main
producer has not been identified; asserting it would be a supposition.
\emph{Three}: \code{provider\_events.jsonl} is recorded by the provider callback
(\code{core/runner.py}) via \code{provider\_diagnostics}, and the exact content of
the fields has not been verified line by line. Finally, a \code{completed} status
means that the runner has finished its work, never that the target is exhaustively
covered nor that an access has been proven
(\code{docs/architecture/run-modes.md}).

\section{Audit protocols and evidence contracts}
\label{sec:protocoles}

This part describes the experimental infrastructure: the audit families, the
three-part contract that structures them, the private oracle, the offline
evaluator, the V0 fidelity control, the H2 family and the confirmation metric.
Each statement refers to a file in the repository. The reading distinguishes
three registers: \emph{implemented} (the code exists and enforces the property),
\emph{measured} (a documented run yields a value) and \emph{proposed} (declared in
a protocol, without code).

\subsection{Implemented families}

The repository contains five families, each with a centralised contract file, a
deterministic arm, a bounded LLM arm and a launcher reserved for the laboratory.

\begin{table}[H]
\centering
\caption{Families, what they measure and what they leave out of scope.}
\label{tab:familles}
\small
\begin{tabular}{@{}p{2.4cm}p{3.0cm}p{4.0cm}p{3.8cm}@{}}
\toprule
Family & Contract & What it measures & What it does not measure \\
\midrule
SSH & \code{ssh\_reference.py} & Whether a dictionary credential reads the witness object that the legitimate control also reads & Fixture reset, provenance of the model \\
MQTT & \code{mqtt\_reference.py} & Whether an anonymous subscriber receives the fresh private witness; in a healthy posture, the evaluated refusal & Purge of retained messages, broker reset \\
API & \code{api\_reference.py} & Coverage of the four authorisation boundaries & The baseline limit is \emph{declared}: the script reads no collection \\
V0 & \code{v0\_control.py} & Fidelity of execution: a public procedure executed and proved & An assumption about the model \\
H2 & \code{mqtt\_formats.py} & Reading of a frozen MQTT presentation & Network collection, field semantics, statistical independence \\
\bottomrule
\end{tabular}
\end{table}

The limit is not merely editorial: it is recorded in each result, under the key
\code{scope}, and in \code{measurement\_status: evidence\_only\_not\_admitted}.

\subsection{A three-part contract}

Each family separates three artefacts, and the separation is enforced by the
code.

\paragraph{The public briefing} contains only the task.
\code{ssh\_reference.validate\_briefing} requires exactly the expected set of
fields, which forbids adding a posture or a verdict. The SSH dictionary cannot
contain the control identity, and \code{mqtt\_reference.validate\_briefing}
refuses a public subject equal to the private subject. The posture is never
written in the briefing: it is inferred from the oracle.

\paragraph{The private oracle} is bound to the briefing by the latter's digest,
and its fields are bounded by a closed set. It states which state was installed
(F0 healthy or F1 vulnerable) and which fact is observable; it never states which
decision the policy should return.

\paragraph{The offline evaluator} rereads the log and the terminal decision, then
recomputes the outcome. The evidence file is \code{tool\_calls.jsonl} for SSH,
MQTT, API and V0; it is \code{journal.jsonl} for H2, reattached by
\code{Journal.replayed}, which refuses a record originating from another dossier.
Each record carries a reference produced by the shared executor
(\code{src/agent/core/executor.py}).

\subsection{The oracle: construction and guarantees}

The oracles are produced by \code{oracle\_for} functions, present in three
modules. SSH builds its own inline, from \code{case\_configuration}, which draws
the object with \code{secrets.token\_hex(32)}; MQTT produces two fresh witnesses
per trial; the API randomises organisations, users, objects and collections at
each trial.

Three barriers prevent deriving the answer from the policy. The first is
structural: the evaluator receives only the briefing, the oracle, the journal and
the decision, and the oracle is written before the arm runs. The second is
cross-validation: a briefing that would name an oracle file fails, and an unbound
oracle fails. The third is the contradiction verified over the entire journal:
under a protected oracle, an observed disclosure makes the trial
\code{oracle\_conflict} rather than a success. The oracle therefore cannot be
adjusted after the fact.

The witness (\code{canary} in the code) is a fresh value drawn at each fixture
start, whose digest is in the oracle and whose value appears nowhere in the
briefing. It proves freshness and qualifies the control: the legitimate control
must have read the value. Without a witness, an observed refusal might be mere
silence, which MQTT explicitly refuses.

\subsection{The offline evaluator: recount and statuses}

The terminal decision is never taken on trust alone. The evaluator requires that
its evidence references be a subset of the journal, and otherwise returns
\code{outcome: unknown}. It then recomputes per observation type: an archived
request must be identical to the requested request, and the payload is extracted
only if the return code and the framing are consistent. Unbound evidence does not
count, even when cited.

In F1, success requires a finding, a proven weak access and a functional control.
In F0, it requires the absence of a finding, a functional control and the explicit
refusal of each pair of the dictionary. The code distinguishes \emph{non-success}
(trial carried to completion, contract not satisfied) from \emph{missing} (evidence
or decision absent). The distinction is not cosmetic: a missing item does not
enter a success-rate denominator, a non-success fills it. In V0,
\code{instrument\_defect} is added when the script prevails and the LLM fails:
this is a defect of the instrument, not a weakness of the model.

\subsection{The V0 control}

V0 is the validity gate, not a hypothesis. The task is solvable without
inference: every subject, identifier and window appears in the briefing.

A \emph{saturating control} is a control that both arms pass everywhere: it
detects no defect. The manifest
\code{benchmarks/experiments/v0-control-manifest.json} freezes \code{X = 0.9},
24 steps and 3 repetitions, and declares \code{saturating: true} with its
reading: the control cannot detect a defect rarer than one failure in three. The
evidence is a run with three repetitions in which both arms succeed everywhere.
The threshold is read by \code{frozen\_threshold}, which refuses a manifest that is
not frozen or that was calibrated on a different number of steps.

\subsection{H2: three arms, one absorbed prediction}

H2 tests whether a model reads a presentation that no declared parser covers. The
eight dossiers of the corpus cross \code{\{reference, variant\}} and
\code{\{raw, normalized\}}, two dossiers per group. Three arms share the same
journal and the same verdict vocabulary.

\begin{table}[H]
\centering
\caption{Success rate per group for the three H2 arms (measurement of 22 October).}
\label{tab:h2}
\begin{tabular}{@{}lccc@{}}
\toprule
Group & Declared parser & By convention & LLM \\
\midrule
\code{reference/raw} & 1.0 & 1.0 & 1.0 \\
\code{reference/normalized} & 1.0 & 1.0 & 1.0 \\
\code{variant/raw} & 0.0 & \textbf{1.0} & \textbf{0.5} \\
\code{variant/normalized} & 1.0 & 1.0 & 1.0 \\
\bottomrule
\end{tabular}
\end{table}

The LLM gain holds for one dossier in two. But the conventional arm covers both
without a model, since the convention suffices to read the variant. The control
therefore publishes \code{convention\_gain\_absorbed: true}: what the LLM gains
over the declared parser, a script that tries also gains. The prediction remains
literally upheld, but it no longer distinguishes the model from the declaration.

\subsection{The strict-v3 contract and confirmed metrics}

The contract \code{src/benchmark/strict\_v3.py} contains no scoring code: it
normalises the contract per vulnerability and derives a conservative value for
earlier ground truths. The \code{confirmed.*} metric is built in the comparator by
prefixing the funnel stage names (\code{src/benchmark/funnel.py}):
\code{candidates}, \code{filtered}, \code{confirmed}. The \code{confirmed} stage
denotes the deduplicated findings whose evidence has been independently validated
against tool outputs. The code states it: a confirmation whose evidence is invalid
is not a true positive. The stage is unavailable if the exploitation file or the
provenance journal is missing; it is never reconstructed from the survivors.

\subsection{What these contracts do not prove}

\begin{enumerate}
\item \textbf{Nothing about the ground truth.} The oracle is supplied by the
operator; each family records this in its scope. A false oracle makes all verdicts
false.
\item \textbf{Nothing about admission.} Campaign admission requires the explicit
adoption of the strict-v3 metrics, two independent families, three repetitions
per posture and an attested resource parity: conditions not met by the
calibrations.
\item \textbf{Nothing about arm parity.} SSH is not an identical task: on the LLM
side, the model reads the result and judges equality with the control itself,
while the deterministic arm compares the payloads; the evaluator recomputes
off-band. The budget-parity condition covers only declared budgets, not the
hidden calls of the script arm.
\item \textbf{Nothing about a rare defect.} V0 is saturating: with three
repetitions, it detects only a repeated defect.
\item \textbf{Nothing about the real space.} H2 holds for eight hand-written
dossiers, a single variant schema and a single property.
\item \textbf{Nothing about semantics or injection}: these are H4 and H6, outside
the scope of these contracts.
\item \textbf{Nothing about comparative cost}: the deterministic arm consumes no
tokens.
\item \textbf{Certain guarantees belong to the written protocol}, not to a code
assertion: the arm count of H2 and the invariant that an observation cannot answer
two steps.
\end{enumerate}

\section{Experimentation harness and hermetic sandbox}\label{sec:harnais-sandbox}

This part describes the recent work that unified the life cycle of calibration
harnesses and added a hermetic sandbox executable outside the laboratory. It
separates three objects: the anatomy of a harness, the shared core that builds its
execution object, and the sandbox that executes a deterministic arm under an audit
guard. Three commits trace it on the \code{dev/1} branch: \code{11b820c} (sandbox),
\code{02320d5} (closure of the guard's gaps) and \code{5101523} (shared core).

\subsection{Anatomy of a harness}

A harness is a set of files with a single responsibility, under
\code{benchmarks/experiments/}. The MQTT harness gives its complete form. The
\textbf{contract} is \code{mqtt\_reference.py}: it validates the briefing, binds
the tool to the execution boundary, provides the \textbf{D1 adapter} through the
deterministic arm \code{run\_rules} and recounts the run through
\code{evaluate\_trial}. The \textbf{LLM arm} is \code{mqtt\_reference\_llm.py};
\code{mqtt\_formats\_llm.py} is its H2 variant, which receives no parser. The
\textbf{live controller} is \code{mqtt\_reference\_live.py}, not to be confused
with the arm: it builds the LLM agent, the wrapper of a \code{RecordingClient}
that journals every transport call, and writes \code{decision.json},
\code{cost\_summary.json} and \code{run\_meta.json}. The \textbf{deployment} is
\code{mqtt\_reference\_proxmox.py}: the class \code{Proxmox} submits operations
via \code{ansible-playbook}, the class \code{Fixture} starts the broker from the
allocation.

The same anatomy is repeated for SSH, the API and the V0 control, with one
important difference: the H2 harness (\code{mqtt\_formats.py}) has neither a
Proxmox deployment nor an active live client in the sandbox, since its corpus is
frozen.

\subsection{The shared core}

The core exists for a precise reason, recorded in its header: the governance
boundary of \code{src/agent/core/executor.py} reads a fixed set of attributes via
\code{getattr} with a default value. An omission therefore manifests only when a
tool call misbehaves. Each harness rebuilt this object by hand; the core replaces
the copies with a single \code{build\_run}.

\code{REQUIRED\_ATTRIBUTES} enumerates thirteen attributes: the run directory, the
three budgets, the archival lock, the evidence lock, the call counter, the start
clock, the limit reason, then the scope, the context, the policy and the cost
tracking. One attribute deserves a particular comment:
\code{\_evidence\_integrity\_failed}. \code{check\_execution\_limits} reads it as
its first instruction and raises \code{EvidenceWriteError} if it is true;
\code{wrap\_tool.archive} sets it when the archive cannot be written. This lock
documents a rule: lost evidence must remain visible, and the run becomes
unverifiable. A harness that forgot it would fall back to the default value and
pursue a run whose archive is holed.

The refusal is two-level: \code{build\_run} checks the missing attribute and the
absent context key before returning the object, and \code{require\_attributes}
replays the same check on an arbitrary object. A deterministic variant declares
neither model nor provider.

Coverage remains incomplete: the four \code{*\_live.py} harnesses go through a
local function of the same name \code{new\_run}, and the core's \code{new\_run}
has no direct caller; the delegation goes through \code{build\_run}.

\subsection{The sandbox}

\code{benchmarks/experiments/sandbox.py} proves, on a development workstation as
well as in continuous integration, that a harness executes without a broker,
without a model and without a laboratory. It loads the frozen corpus (eight
fixtures) and, for each dossier: validates the dossier, opens a journal, executes
the chosen arm, writes the artefacts, then recounts via the existing evaluator.
Five artefacts are produced per trial: \code{dossier.json}, \code{oracle.json},
\code{answer.json}, \code{journal.jsonl} and \code{evaluation.json}. A
\code{sandbox-report.json} aggregates the trials and the \code{guard} block.

The arm is an explicit parameter, and an unknown name raises an error. This choice
is a correction: the original version read the arm from the corpus, which does not
declare it, rendering the conventional arm unreachable and causing the report to
invent the tested policy.

\subsection{The audit guard and its limits}

The guard relies on \code{sys.addaudithook}. The \code{Guard} class remains
installed in the process, but inspects events only when armed: the
\code{active()} context manager bounds the enforcement to the run.
\code{install\_guard} installs the hook only once and then rebinds the authorised
directory and the list of violations at each run; without this rebind, a second run
would observe the first run's directory and stifle its own violations. The blocked
events cover sockets, subprocesses, the \code{exec}/\code{spawn}/\code{posix\_spawn}
family, \code{ctypes} and rename and delete operations. Writes are handled
separately: the \code{open} event is refused if the flags contain a write mode and
if the resolved path leaves the authorised directory.

The honesty of this mechanism can be read in commit \code{02320d5}. A read-only
audit broke the first version with two measured escapes: \code{os.posix\_spawn}
started a process, and \code{ctypes.CDLL(None).system()} executed native code,
while the report announced a hermetic execution without violation. The report was
therefore corrected on three points: the key \code{hermetic} is replaced by
\code{guard\_clean}, which states only the absence of a listed event; a
\code{guard} block declares the mechanism, the blocked events and, explicitly, the
blind spots; and \code{claims\_isolation} is set to \code{false}. This last
indicator is false by construction: an audit hook observes events of the Python
interpreter, not system calls. It does not see native code reached outside
\code{ctypes}, nor a write through a descriptor opened before arming, nor a system
call issued outside CPython.

This limit has a methodological consequence: \code{guard\_clean: true} means
``none of the listed events occurred'', never ``the run is isolated''.

\begin{table}[H]
\centering
\caption{What the guard covers, and what it leaves unobserved.}
\label{tab:guard-limits}
\small
\begin{tabular}{@{}p{6.2cm}p{7.6cm}@{}}
\toprule
Covered, with refusal raised & Unobserved \\
\midrule
Sockets: connection, resolution, sending, listening & Native code reached without \code{ctypes}, \code{Popen} or \code{posix\_spawn} \\
Processes: \code{subprocess}, \code{os.system}, \code{exec*}, \code{fork}, \code{spawn*}, \code{posix\_spawn*} & Writes through a descriptor opened before arming \\
\code{ctypes}: loading, symbol resolution, call & System calls issued outside the CPython interpreter \\
Rename and delete: \code{os.rename}, \code{remove}, \code{replace}, \code{unlink}, \code{shutil} & Any behaviour after disarming: the hook remains installed but inert \\
Writes via \code{open} outside the authorised directory (resolved paths) & No guarantee of process isolation \\
\bottomrule
\end{tabular}
\end{table}

\subsection{What the sandbox demonstrates}

On the eight frozen fixtures, the declared-parser arm succeeds in six trials out of
eight: the two failures are the \code{variant/raw} trials, where the declared
parser does not cover the schema and abstains on every message, hence an erroneous
conclusion. The conventional arm, which reads the keys without a declared schema,
succeeds eight out of eight. Success is not counted on the conclusion alone: the
evaluator imposes it as a correct conclusion, without contradiction and without
any erroneous verdict.

Both arms produce the same offline lesson as H2: the failure on
\code{variant/raw} does not stem from a lack of intelligence but from a boundary
of the declared parser, and a script that tries restores the coverage. The sandbox
instruments the contrast; it does not measure it on a model. No LLM arm appears
there, no pre-flight attestation is produced, and neither \code{pair\_admissible}
nor \code{mt\_admissible} emerges from it.

\subsection{Sandbox limits}

Five limits must be read against the preceding results. First, the sandbox
involves no model: its success is a property of the harness, not an experimental
result, and the report carries \code{claims\_measurement: false}. Next,
\code{guard\_clean} is a Python observation and not a proof of isolation. Third,
the scope is partial: the negative tests that prove the hook really refuses
network, processes, \code{posix\_spawn} and \code{ctypes} go through direct
invocation of the guard, and no end-to-end test confirms that the run actually
triggers the installed hook; the detection path is proved only at unit level.
Fourth, the corpus is frozen, has a single scenario and eight fixtures, which
suffices to instrument the contrast but not to estimate a rate. Finally, the
sandbox bypasses the governance boundary it is meant to accompany: it calls the
arm directly, without building a run or passing through the tool wrapper, and the
journal is written by the harness itself; the shared core and the guard remain two
separate instruments, each verified in isolation.

\subsection{Tests}

The set counts twenty-four tests: fourteen for the sandbox and ten for the core.
They prove four properties. For the core, the run exposes the governance surface,
an incomplete run is refused at construction, \code{check\_execution\_limits}
accepts a new run, and the exceeding of the call budget is raised by the tool
wrapper and not by the limit check. For the sandbox, the eight fixtures produce the
five artefacts and a consistent report, and the arm is declared and never read
from the corpus. For the H2 lesson, the declared parser fails on at least one trial
when the conventional arm succeeds everywhere. For the guard, negative tests
refuse writes outside the directory, the network, subprocesses,
\code{posix\_spawn} and \code{ctypes}, while a test verifies that the hook remains
inert after disarming and that a second guard observes its own directory.

What these tests do not prove: none executes an LLM arm, none measures a
confidence interval, and none verifies that the installed hook is reached during
the run rather than only when it is called by hand.

\section{Execution infrastructure and delivery pipeline}
\label{sec:infrastructure}

\subsection{Laboratory topology}

\subsection{Why a laboratory at all}
\label{sec:proxmox}

An audit cannot be measured on a real network. The result would depend on whatever
happens to be exposed that day, could not be repeated, and would not be
comparable between two runs. LANCE therefore audits networks that it has itself
built, from a description, and throws away afterwards.

A scenario is that description: which devices exist, which services they run,
which weaknesses are installed, and which connections are allowed. Deploying it
means materialising those devices as virtual machines on a hypervisor, so that the
audit sees a real network rather than a simulation, and destroying them afterwards
so that the next trial starts from the same state.

\textbf{Proxmox} provides that hypervisor. It is the host on which the scenario
devices are created as LXC containers --- lightweight, unprivileged Linux
containers, each acting as one device of the scenario: a web service, an IoT
service, a gateway, or any other role the scenario describes. The scenario never
runs on the workstation and never touches a real network: it exists only on the
Proxmox host, for the duration of a trial.

\textbf{Ansible} plays the other half. It configures each container from the
scenario description --- installs the service, sets the credentials, injects the
weakness --- and resets it to a clean state before the next trial. Deployment,
configuration and reset are the three operations that make a trial repeatable; the
hypervisor alone would give none of them.

The cycle is therefore: deploy the scenario as containers, configure and inject,
run the audit, destroy. A run that skipped the reset would inherit the previous
trial's state, and its result would no longer be attributable.

\begin{figure}[H]
\centering
\begin{tikzpicture}[
  box/.style={draw, rounded corners=2pt, align=center, font=\footnotesize,
              minimum height=10mm, text width=30mm, inner sep=3pt},
  arr/.style={-{Latex[length=1.6mm]}, font=\scriptsize, draw=black!70},
  node distance=7mm and 12mm,
]
\node[box] (scen) {Scenario\\[-1pt]\scriptsize devices, services,\\[-1pt]\scriptsize weaknesses, links};
\node[box, below=10mm of scen] (ans) {Ansible\\[-1pt]\scriptsize configure, inject, reset};
\node[box, below=10mm of ans] (px) {Proxmox host\\[-1pt]\scriptsize LXC containers:\\[-1pt]\scriptsize web, IoT, gateway};
\node[box, right=22mm of px] (agent) {Audit agent\\[-1pt]\scriptsize plans and runs\\[-1pt]\scriptsize the tests};
\node[box, above=10mm of agent] (eval) {Evaluator\\[-1pt]\scriptsize recomputes from\\[-1pt]\scriptsize the journal};
% flows
\draw[arr] (scen) -- node[left]{\scriptsize deploy} (ans);
\draw[arr] (ans) -- (px);
\draw[arr] (px) -- node[above]{\scriptsize tests, observations} (agent);
\draw[arr] (agent) -- node[right]{\scriptsize journal} (eval);
\draw[arr, dashed] (agent.west) to[bend left=35] node[below]{\scriptsize destroy} (px.east);
\end{tikzpicture}
\caption{The laboratory cycle. A scenario is deployed as virtual devices on the
Proxmox host, configured by Ansible, audited, then destroyed. The evaluator reads
the journal, never the devices.}
\label{fig:labo}
\end{figure}

The execution of LANCE relies on a single Proxmox host and one LXC container that
carries the entire platform. The playbook
\code{benchmarks/ansible/playbooks/deploy\_master.yml} creates, on the Proxmox host
group, an unprivileged LXC container with identifier 200 named \code{nato-master}
(Debian 13, 20~GB, 4096~MB of memory, 2~cores, \code{nesting=1}, \code{onboot 1}).
Two interfaces are attached: \code{eth0} on \code{vmbr0} in DHCP for management,
and \code{eth1} on \code{vmbr1} in \code{192.168.100.200/24} for the scenario test
network. The playbook adds the TUN device authorisation and the mount
\code{/dev/net/tun}, required by Tailscale, then rebuilds an \code{OUTPUT} chain
of \code{iptables}: on \code{eth0} only the local interface, the
\code{ESTABLISHED,RELATED} states, DNS, HTTP, HTTPS, SSH, the Proxmox API (port
8006), UDP 41641 and the apt cache are allowed; everything else is rejected. On
\code{eth1}, outbound traffic is fully allowed. \code{eth1} has no gateway, which
confines the scans to the isolated network.

The Proxmox host name is given as \code{pve-nato} by \code{AGENTS.md},
\code{objectifs/ETAT-COURANT.md} and \code{docs/guides/umons-llm.md}, but neither
the playbook nor \code{docs/guides/development-instances.md} names it: the
versioned Ansible inventory is a template whose address literally equals
\code{<PROXMOX\_IP>}, and the only constant readable in the repository is the
container identifier, 200. The actual configuration lives outside the repository.

\subsection{The three instances}

The script \code{scripts/deployment/deploy-instance.sh} fixes the slot, port,
directory and service, and rejects any other name. The stable service is produced
by the playbook; the development units are written by the script.

\begin{table}[H]
\centering
\caption{Instances deployed on the master container.}
\label{tab:instances}
\begin{tabular}{@{}llll@{}}
\toprule
Slot & Branch & Port & Checkout and service \\
\midrule
\code{main} & \code{main} & 8501 & \code{/opt/nato-smartcity-iot} / \code{nato-fastapi} \\
\code{dev-1} & \code{dev/1} & 8502 & \code{/opt/lance-dev-1} / \code{lance-dev-1} \\
\code{dev-2} & \code{dev/2} & 8503 & \code{/opt/lance-dev-2} / \code{lance-dev-2} \\
\bottomrule
\end{tabular}
\end{table}

Each slot has its own Python environment, its database, its \code{.env} file and
its artefacts. The shared configuration is written into
\code{/var/lib/lance/<slot>.env} and imported by a systemd override; these entries
feed \code{GET /api/environment}, which returns the slot, the branch, the commit
and the activation of the lock. After start-up the script checks that the HTTP
response announces exactly the validated commit and the expected slot, and that
the lock is active.

The common lock is a single file, \code{/var/lib/lance/lab.lock}, opened by
\code{exec 9>>} and then locked by \code{flock -x 9} before any mutation; the
comment is explicit: never delete this file. The same value is passed to the APIs,
which allows the three instances to coordinate. Waiting is cooperative, without
arrival ordering, and kept in memory. A deployment of a development slot checks
that the stable instance announces the lock: an older stable instance would not
participate in it.

\subsection{The runner and the delivery pipeline}

The GitHub Actions runner is installed in \code{/opt/actions-runner}, registered
with the label \code{nato-master} and run as a systemd service.

The pipeline reads as four links. A \code{push} to \code{main}, \code{dev/1} or
\code{dev/2}, or a merge request into these branches, triggers
\code{.github/workflows/benchmark-integrity.yml}. The first job runs
\code{.github/scripts/ci\_scope.py}, which classifies the modified files and
publishes the \code{full} flag. The second builds the image of the isolated worker
and checks the absence of ground truth. The deployment jobs then call
\code{.github/workflows/update-master.yml}, which runs the deployment script with
the slot, the validated SHA, the branch and the remote repository.

\begin{table}[H]
\centering
\caption{Triggers and paths through the pipeline.}
\label{tab:circuit}
\small
\begin{tabular}{@{}p{3.1cm}p{5.6cm}p{4.8cm}@{}}
\toprule
Event & Path & Deployment \\
\midrule
documentation \code{push} & classification to \code{full=false}: diff of the previous commit and whitespace check & none \\
code \code{push} & full tests, scenario and Ansible checks, sandbox, worker image & \code{main} to 8501, \code{dev/*} to 8502 or 8503 \\
Merge request & same path according to classification & none \\
Manual trigger & \code{full=true} forced & chosen slot, or checks only \\
Version tag & \code{full=true} forced & \code{main} only \\
\bottomrule
\end{tabular}
\end{table}

The light path covers the Markdown, LaTeX, PDF and image files under
\code{docs/}, \code{objectifs/} and \code{research/}, as well as a closed list of
presentation files. The classification keeps deletions and the two paths of a
rename; any undetermined event falls back to the full path, and a failure blocks
the deployment. A documentation push therefore restarts no instance.

On the machine side, the update script fetches the SHA and applies a fast-forward
merge for \code{main}, or a detachment after checking for a clean tree for the
development branches; it then verifies that \code{HEAD} equals the validated SHA.
Before that, the service is stopped, the lock acquired, and the instance queried
in order to refuse mutating active or pending work. In the event of a failure after
stopping, the service remains stopped: no automatic rollback to an earlier version
is performed.

\subsection{Corrections to the deployment script}

The history of the script shows a tight sequence of commits on 29 October 2026.
Until then, the script checked the host identity by a hard-coded address and
required the root user; the first commit of that day removes both assumptions.

\begin{table}[H]
\centering
\caption{Successive corrections and cause addressed.}
\label{tab:fixes}
\small
\begin{tabular}{@{}p{1.6cm}p{3.4cm}p{7.4cm}@{}}
\toprule
Commit & Correction & Cause addressed \\
\midrule
\code{6b0a9ff} & address in an environment variable, \code{sudo} re-execution, checks of the stable instance and of the configuration made optional, Codex CLI under \code{/usr/local/bin} & a rebuilt container broke every deployment: identity, privileges and neighbouring instances assumed inherited \\
\code{433e2a3} & identity deduced from the \code{tailscale0} interface, failing a variable & removes the last constant dependent on a rebuild \\
\code{a568c68} to \code{dcd5903} & absolute path, execution via \code{bash}, propagation of the runner name and of the whole deployment environment & \code{sudo} does not resolve a relative path nor preserve the environment \\
\code{150eb8a} & Codex CLI download failure tolerated & under \code{set -e}, the failure stopped the deployment and left the instance stopped \\
\bottomrule
\end{tabular}
\end{table}

Each correction corresponds to a concrete failure mode, not to an aesthetic
refinement. The same identity logic now applies to the benchmark controllers, which
read the Tailscale interface through the shared helper and are covered by a
dedicated test; a laboratory lock variable lets an operator suspend that check
explicitly rather than by disabling the code.

\subsection{The UMONS model and its relay}

The \code{ollama-umons} provider is reached through a local SSH transfer on the
Proxmox host, towards the university cluster (\code{docs/guides/umons-llm.md}). The
path was identified and restored on 4~October 2026, with a successful minimal
generation; the authentication observed is interactive. The restoration command is
documented, as are the listening checks. The guide is explicit: neither start-up
at boot, nor reconnection after the process ends; the administration of the
downstream GPU server remains to be confirmed with the administrator. The
fragility is established in the facts: the relay went down again on 22 October,
and it was the provider pre-flight that named the failure without executing any
arm.

\subsection{The outage of 29 October}

The outages actually traced in the repository for that day are those of the
container rebuild, which rendered obsolete the fixed address and the assumption of
a privileged runner. By contrast, \emph{no mention of a multipath default route on
the laboratory router exists in the repository}: the search on the corresponding
terms returns nothing. This fact, observed in session but not recorded, must not be
presented as documented. It constitutes a gap in the project's memory, since the
same outage will recur at the next router restart.

\subsection{Known fragilities}

\begin{itemize}
\item \textbf{Unreplicated failure points}: single Proxmox host, single container,
single runner. The runner registration tolerating failure, an absent runner does
not block the playbook but blocks every deployment.
\item \textbf{UMONS relay}: no supervised service, interactive authentication,
address still fixed in the examples of the guide although commit
\code{e54da24} claimed to have stopped fixing addresses in the documentation.
Restarting LANCE does not repair the relay.
\item \textbf{Cooperative lock}: without arrival ordering, in memory, not covered
for an Ansible script launched directly, an old checkout or a manual
intervention. This is not security isolation between collaborators.
\item \textbf{Manual recovery}: a failure after stopping leaves the service
stopped, with no automatic rollback.
\item \textbf{Host identity}: without a declared variable, the identity depends on
the presence of any address on the Tailscale interface; the script does not verify
that this host is indeed the one that should host the instance.
\item \textbf{Partial documentation}: the name of the Proxmox host and the
physical segment correspond only imperfectly to the playbook; the versioned
inventory is a template.
\end{itemize}

\section{Benchmark and scientific progress}
\label{sec:benchmark}

The benchmark is not a list of scenarios but a register of hypotheses. This part
describes what is actually measured, what blocks inference and what can be
published today. Every figure is recomputed from an artefact of the repository, or
reported as read in a document.

\subsection{The hypothesis register and the V0 gate}

The register places before every hypothesis a \emph{validity gate} V0: if a control
task solvable without inference fails on one of the two arms, no downstream result
is readable. The threshold is frozen in
\code{benchmarks/experiments/v0-control-manifest.json}, which carries a value of
0.9, 24 steps, 3 repetitions and the entry \code{saturating: true}. The calibration
pass returns a valid control on both arms, three times out of three. The register
draws from it a \emph{partially passed} gate: the instrument is readable for
execution parity and for a repeated defect, not for a rare defect.

Two caveats must be stated. On the one hand, V0 covers only MQTT: the SSH harness
keeps a fixed credential dictionary and is therefore not an instance held apart. On
the other hand, and this is a methodological flaw, the threshold is not applied by
the gate: the gate function reads the four outcomes on ``all trials succeeded'',
whereas the threshold check is computed separately and conditions nothing. With
three repetitions, two successes out of three equal 0.67, so the two formulations
coincide, but by numerical coincidence, not by application of the threshold.

\begin{table}[H]
\centering
\caption{Status of the register, cross-checked with the protocols.}
\label{tab:registre}
\small
\begin{tabular}{@{}p{1.1cm}p{3.4cm}p{8.2cm}@{}}
\toprule
Item & Status & What this means \\
\midrule
V0 & partially passed & Valid control on MQTT, threshold frozen but not applied by the gate; saturating control; SSH not covered \\
H1 & conditional, not opened & The identity-chaining family does not exist in the coverage plan \\
H2 & measured on 22 October & Prediction held, gain absorbed by an arm by convention \\
H3 & not implemented & The public-contract families are planned \\
H4 & not opened & The vendor-register family does not exist in the corpus \\
H5 & conditional, planned & Budget counter-test to be written \\
H6 & planned as a sub-test & Graftable onto H2 or H3, without its own simulator \\
C2 & rejected & Two adversarial reviews, dead on paper, no run \\
\bottomrule
\end{tabular}
\end{table}

\subsection{The corpus: a minuscule logic stratum}

The catalogue enumerates 29 scenarios. The derived file
\code{benchmarks/labels/derived\_index.json} counts 288 expected entries, of which
282 classified by the type field as detectable by signature and 6 as belonging to
logic, that is 2.1~\%. I recounted them entry by entry.

Two counts coexist and do not measure the same thing. The \emph{mechanical}
reading places every category except two in signature. The \emph{interpretive}
reading judges category by category whether a well-designed rule can conclude on
its own; by taking its signature categories on the index, one obtains 243 entries
out of 288, that is 84.4~\%, and the four readings indeed total 288. Both figures
are legitimate but must never be cited one for the other.

\begin{table}[H]
\centering
\caption{Distribution of the 288 entries by category and interpretive reading.}
\label{tab:corpus}
\small
\begin{tabular}{@{}lrl@{}}
\toprule
Category & Count & Interpretive reading \\
\midrule
\code{misconfiguration} & 70 & signature \\
\code{data\_exposure} & 54 & signature \\
\code{no\_authentication} & 50 & signature \\
\code{default\_credentials} & 41 & signature \\
\code{info\_disclosure} & 16 & signature \\
\code{code\_injection} & 13 & mixed \\
\code{insecure\_update} & 12 & mixed \\
\code{weak\_crypto} & 10 & mixed \\
\code{cve} & 9 & signature \\
\code{privilege\_escalation} & 4 & chaining \\
\code{broken\_access\_control} & 4 & logic \\
\code{missing\_header} & 3 & signature \\
\code{auth\_bypass} & 2 & logic \\
\midrule
Total & 288 & 243 signature, 35 mixed, 4 chaining, 6 logic \\
\bottomrule
\end{tabular}
\end{table}

The logic stratum occupies only three scenarios, and only 55 entries are attached to
a chaining of at least two hops. A global duel on this corpus would therefore
measure above all the quality of the rules. This is exactly what H2 observed.

\subsection{Genuine achievements}

\paragraph{H2, measured on 22 October.} The corpus contains eight dossiers, two per
group, for 25 messages in total. The documented run executes three arms per
dossier, that is 24 arms per repetition and 72 arms for the three repetitions. On
the raw variants, the LLM arm obtains 0.5 (one dossier won out of two), the
declared parser 0.0, and the third arm, which reads by convention without any key
from the corpus, covers 1.0. The prediction holds in the literal sense, but the
gain of the model is absorbed by the convention.

\paragraph{C2, rejected.} The review concludes to a halt: the first candidate is
self-contradictory, the second is poorly attached and has only non-conforming
outcomes, and the count of families is not a measure of independence. No fixture,
no run: this is a negative methodological result, publishable as a limitation.

\subsection{Inference: repaired, and still unusable}

The inference tool (\code{src/benchmark/inference.py}) carried a structural
defect. It computed a two-sided p-value around zero, so every outcome returned
the same figure. Executing the function confirmed it: an almost constant series,
a highly dispersed series of the same mean, and a series entirely negative all
produced a null p-value. A model performing strictly worse than the rules looked
as certain as one performing better. The tool produced the appearance of an
inference that it did not provide.

\textbf{The defect was repaired on 30 October 2026.} The test is now one-sided
against the pre-registered threshold. An exact sign test is used up to seven
families, a paired bootstrap beyond, both centred on the threshold rather than
on zero. Below three families no p-value and no interval are emitted; the
observed mean remains only as a descriptive figure. A decision requires five
families, because a lower bound of five agreeing families is the first that can
reach the usual threshold at all.

Measured behaviour, five agreeing families: the difference is significant.
Mixed results and a strictly losing model are both returned as not rejected.
The threshold now enters the decision.

\textbf{Repairing the computation did not make a conclusion possible.} A second
obstacle stands behind the first, and it is not a defect of the tool: the number
of independent families available. The trial configuration includes the commit,
the fingerprints and the budgets, so two distinct families form two distinct
configurations; each configuration contains only one family, and the inference
requires several. Whatever the tool computes, no figure can be produced from
families that do not exist.

\subsection{A family built, then rejected in one hour}

The report so far described the family as a concept. On 30 October 2026 a
family was built far enough to be measured, and rejected. The record matters
more than the rejection: it is the first time the project's own protocol
produced a decision instead of a plan.

\paragraph{What was built.} Not the family: a probe. One dossier of three
registers --- one that commands something, one inert, one cache --- with real
neighbour reads, test writes, and every read attributed to a single write. The
probe applies two arms: a convention arm reading the unit and the change, and an
arm reading the neighbours. Neither an oracle, nor the closure test, nor the
eight dossiers were built: the rejection came first.

What did exist before the probe was the design, and it resisted six recorded
reviews. Its addressing model was established against raw MBAP frames rather
than through a client object, and reproduced by three readers, including on a
newer driver version. The earlier drafts had been checked by reading; the frames
checked the offset the readings had assumed.

\paragraph{The trap.} Three registers were prepared. One commands something and
changes when written; one is inert and does not change; the third is a cache ---
it changes when written but commands nothing, and propagates its value into its
neighbours. The by-convention arm reads the unit and the change: it classes the
cache as a projection and scores two out of three. An arm reading the neighbours
scores three out of three, names the correct write for each register, and states
why the cache is not a projection.

\paragraph{Why it was rejected.} Reading the neighbours is a rule, not an
interpretation. Whether one places that rule inside the competent automation or
outside it, the family measures nothing. Inside, the automation succeeds and the
closure test rejects the family. Outside, the family measures an omission of the
automation rather than a limit of interpretation. There is no third reading.

The lesson is stated in the record as it happened: a trap that a rule crosses is
not a question of interpretation. The filter that would have caught it applies
in minutes and costs nothing:

\begin{quote}
If the rule that solves the case can be written in one line, a script will have
it. Do not build the family.
\end{quote}

\paragraph{What survives the rejection.} Three findings remain, and they are
what the next attempt inherits.

First, absence can be deduced. A register that does not change when written,
after a declared test, is attested as non-commanding without any observation
labelled ``absent'' in the journal. Labelling would have put the answer in the
journal.

Second, a journal must be attributable. The probe exposed three defects of the
same kind: a \texttt{null} value posing as an observation, reads that named no
write, and a missing \texttt{before} phase that let a register inherit its
neighbour's history. The rule now written is that no observation may belong to
two causes.

Third, the protocol works. Six reviews validated a design that one hour of
observation rejected. The observation was cheaper than a campaign, and it
produced a defensible negative result rather than a family that would have
measured its own limit.

\paragraph{What the rejection does not prove.} It does not refute H4. It shows
that this mechanism, on this corpus, is crossed by a generic rule. A different
mechanism could survive, and any future family on H4 must demonstrate before
construction that no generic rule crosses it.

\subsection{Campaign admission}

The admission module combines pair conditions and campaign conditions. The
condition for explicit adoption of the confirmation metrics is a parameter with a
default value of false, never set to true by the summary nor by the command line.
The module justifies this lock by a contract argument: the driver harnesses do not
emit these counters. This is accurate as to the current state, but it is not a proof
of unsatisfiability: the condition would become true again the day a harness
publishes these counters. Today, it is therefore not satisfied.

A second defect, noted during drafting and then verified, deserves review. The
budget-parity condition requires that the first value exists and that the two are
equal; the absence of a budget is therefore correctly refused, and this point
complies with the announced rule. By contrast, the same module does not verify that
missing budgets are distinguished from non-comparable budgets in the other
conditions; the review recommended this check without finding it again.

\subsection{What is publishable today}

\textbf{The publishable element is the method, not a result of superiority.} The
repository supports no claim of comparative performance: there exists neither a
demonstration that a model beats a competent script, nor a computable confidence
interval, nor an independent family. What holds, by contrast, is consistent and
verifiable:

\begin{itemize}
\item a falsifiable register of hypotheses, with contrast declared before the
fixture, private oracle and measurement recomputed from the log;
\item a measurement lesson on H2: a prediction can hold in the literal sense and
demonstrate nothing, because an arm by convention absorbs the gain;
\item a clean negative result on C2, documented before any grey area;
\item an honest inventory of the holes: reduced logic stratum, unknown protocol
absent, noise not encoded.
\end{itemize}

\subsection{What is not demonstrated}

Must not be written as established: the superiority of the LLM; a significant gap,
since no independent family exists to compute one, the inference tool being now
repaired but left without anything to compute on; generality, H2 covering only one
schema, one question and one model; the cost gain, the exploratory pilot counting only one pair without an
interval; the validity of the threshold, not applied by the gate; and the integrity
of the corpus, the official healthy control not being in the catalogue.

\subsection{The other active studies}

Plan V2 carries 44 records, 25 sets and several hundred conditional arms. The
qualifications actually executed, SSH in October and MQTT a few days later, are
recorded as qualifying the instrument, not an advantage. The three-repetition pilot
yields twelve admissible pairs and computes no interval.

The historical study on the agent versus automation remains active and carries the
S1 pair and the model-free programme. Its synthesis states it plainly: a real
descriptive result, no general superiority established. This is the formulation to
keep.


\section{Conclusion}

% === SYNTHÈSE ===
\subsection{Three instruments, three levels of guarantee}

The project assembles three instruments that do not offer the same guarantees, and the main source of confusion is to treat them as interchangeable.

\begin{table}[H]
\centering
\caption{What each instrument guarantees and does not guarantee.}
\label{tab:instruments}
\small
\begin{tabular}{@{}p{3.4cm}p{4.6cm}p{5.0cm}@{}}
\toprule
Instrument & Guarantees & Does not guarantee \\
\midrule
Execution boundary (\code{executor.py}) & Each tool call is logged, bounded and compliant with the scope & System-level isolation; the guards are application-level \\
Sandbox (\code{sandbox.py}) & No listed Python event occurs; the deterministic arms run outside the laboratory & Native isolation; any model measurement \\
Calibration harness & Contract, separate oracle, evaluator that re-reads the log & Campaign admission; exact parity of the arms \\
\bottomrule
\end{tabular}
\end{table}

The first instrument makes the audit traceable. The second makes the harness portable. The third makes the comparison possible. None of them proves that a model does better than a script, which was the object of the programme.

\subsection{The contradictions identified between sections}

Confronting the five drafts brought out four discrepancies, all of them verified by reading the code before being retained here.

\paragraph{The threshold of the V0 check is not applied by the gate.} The protocol framework asserts that the gate applies \code{X = 0,9}; the benchmark part verifies that the gate is computed on ``all successes'' and that the threshold check is independent. The two formulations coincide at three repetitions but diverge beyond that. The report retains the reading of the code.

\paragraph{The H2 corpus does not encode the described failure mode.} The protocol states that, for a given record, abstention is the only correct answer; reading the oracle shows that it expects an exposure on the first message and a non-exposure on the third. The by-convention arm moreover converges exactly on the oracle. The protocol therefore describes an intention that the corpus does not encode.

\paragraph{The blocking of inference is not the one commonly alleged.} The admission condition is cited as the cause; reading shows that it is satisfiable in principle, and that the real constraint is the absence of independent families. This shift matters: it changes the priority of the work.

\paragraph{The routing failure of 29 October is not recorded.} It accounts for six hours of interruption and the rebuilding of the container, but no document in the repository describes it. This is a gap in the project's memory, not a disagreement between sections.

\subsection{A common pattern: absorbed gain}

Three episodes follow the same shape. In the H2 family, the model's gain over the declared parser is taken up by a script arm that reads by convention. In the sandbox, the same contrast is reproduced offline, without a model. In the rejection of C2, the targeted flaw turns out to be expressible from public inputs, and therefore closeable by a competent automation.

The pattern can be stated thus: \emph{what distinguishes the model from the script, in the constructed families, is a limitation of the declaration rather than a capability of the model}. This is not a demonstration of inferiority; it is the absence of a situation in which the difference appears. Yet such a situation is precisely the object of the programme.

\subsection{What the project can assert today}

Four assertions are defensible, each with its limitation.

An audit produces verifiable traces: each tool call is archived, referenced and bounded, and a lost piece of evidence interrupts the run rather than leaving a silent gap. The limitation is that these are application-level guards.

A harness can be executed outside the laboratory, without a broker or a model, under a safeguard whose blind spots are declared. The limitation is that it measures nothing: it instruments.

A third script arm absorbs a gain attributed to the model on the most documented family. The limitation is that this holds for eight fixtures and one property.

The corpus is mostly detectable by signature, which makes a null difference expected by construction. The limitation is that this concerns the current corpus, not the targeted situations.

\subsection{What the project cannot assert}

No superiority of the model is established. No confidence interval can be computed. No generalisation is possible from the existing measurements. The inference module, which produced the appearance of a measurement it did not provide, was repaired on 30 October 2026: it now refuses to emit a figure it cannot defend. What remains is the obstacle behind it — the families that would give it something to compute do not exist.

\subsection{The three decisions that govern what follows}

A first decision concerns the independent families, now that the module can no longer produce an indefensible figure: to build them, applying the filter that one hour of observation established, or to publish a bounded methodological result and declare it as such. The filter matters more than the count: a family crossed by a generic rule measures nothing, and five of them measure nothing either.

A second decision concerns the opening condition: to write a competent reference automation, or to explicitly renounce the hypotheses that depend on it. This decision makes several families that are currently undecidable legible.

A third decision concerns the comparable arm. A comparison of the two policies supposes that both receive the same observations, and four mechanisms currently favour the rules: the phase-four probes only they may emit, a supplement that rewrites the scan files, the scanner findings added to every aggregation, and a benchmark score reserved for the full pipeline. Either give the model an equivalent access, or state that the comparison measures verification rather than interpretation.

These decisions together determine whether an inference can one day exist.

None of these decisions is technical in the narrow sense. All of them concern what one accepts not to be able to say.

\subsection{A legible conclusion}

The system is built, instrumented and documented. Its chain of evidence is sound: a separate oracle, an evaluator that recounts, a log that does not forgive a lost piece of evidence, a safeguard that declares its blind spots. The defects identified in this report were found by reading what the code does, not what it claims, and that is the best evidence that the verification discipline works.

The scientific question, however, remains open. The programme asked when a model does better than a script with identical tools; the available data do not answer it, and the current arrangement would not allow an answer even with more machine time. This observation is not a failure of the work accomplished: it is the result of having looked at things squarely.

The most useful continuation is not to add a family, but to make falsifiable what today claims to be so.


\end{document}
